\section{The cancellation integral and the central model}
\label{app:model-estimates}

This appendix records the integral identities behind the central correction.
All logarithms are natural. Products of a vanishing coefficient with its
logarithm have their continuous values. Integrals below concern complete
cancellation expressions; the separate cotangent integrals generally diverge.

For real $C,D$ and $0<e<\pi$, put
\begin{align}\label{eq:generic-I}
 I(C,D,e)&=\int_0^\pi\bigl\{|C\cot z+D\cot(z+e)|
                  -|C\cot z|-|D\cot(z+e)|\bigr\}\,dz,
 \\
 J(v,s,e)&=I((v-s)/2,(v+s)/2,e).\notag
\end{align}
The integrand has finite one-sided limits at its two pole positions; their
point values do not affect the integral. Changes of variable on the
cotangent circle show invariance under exchanging the two channels,
reversing both coefficient signs, or replacing $e$ by $-e$.

\begin{lemma}[Complete cancellation formulas]\label{lem:generic-cancellation}
Let $v,s\geq0$ and $0<e\leq\pi/2$. If $s\leq v$, set
$c=(v-s)/2$, $b=(v+s)/2$. Then
\begin{equation}\label{eq:J-same-sign}
 J(v,s,e)=2\{c\log(c/v)+b\log(b/v)\}.
\end{equation}
The value is zero when $v=s=0$. If $s>v$, set
$c=(s-v)/2$, $b=(s+v)/2$. When $s\sin e\geq v$,
\begin{equation}\label{eq:J-no-root}
 J(v,s,e)=2s\log(\sin e).
\end{equation}
Otherwise, with $\vartheta=\arcsin(s\sin e/v)$ and
$z_1=(\vartheta-e)/2$,
\begin{equation}\label{eq:J-root}
 J(v,s,e)=2s\log(\sin e)+2\left\{
 b\log\frac{\cos z_1}{\sin(z_1+e)}
 -c\log\frac{\cos(z_1+e)}{\sin z_1}\right\}.
\end{equation}
These formulas join continuously at all their thresholds and zero
coefficients. Writing
\begin{equation}\label{eq:J-R}
 J(v,s,e)=2(s-v)_+\log(\sin e)+R(v,s,e),
\end{equation}
one has
\begin{align}\label{eq:J-R-bounds}
 |R|&\leq\max(2v,s+v)\log2,&
 |R_s|&\leq\log\frac{v+s}{|v-s|},\\
 R_e&=\frac{2\{v\cos e-\sqrt{(v^2-s^2\sin^2e)_+}\}}{\sin e}
       \geq0 &&(s>v),\notag
 \qquad R_e&=0 &&(s\leq v).\notag
\end{align}
The parameter derivative bound is used under integration at $s=v$.
Finally, for
\begin{equation}\label{eq:root-Z}
 E(v,s,e)=J(v,s,e)-2s\log(\sin e),\qquad
 Z(v,s,e)=2\int_e^{\pi/2}
          \frac{\sqrt{(v^2-s^2\sin^2t)_+}}{\sin t}\,dt,
\end{equation}
one has $E\leq Z$ everywhere and $E=Z$ when $s\geq v$.
\end{lemma}

\begin{proof}
On each sign interval the cancellation integrand in
\eqref{eq:generic-I} is $-2$ times the smaller absolute channel. Its
primitive is therefore a linear combination of $\log|\sin z|$ and
$\log|\sin(z+e)|$. For $s\leq v$ the two opposite-sign intervals give
\eqref{eq:J-same-sign}. For $s>v$ the possible additional equality
roots satisfy
$\sin(2z+e)=s\sin e/v$. There are no such roots if
$s\sin e\geq v$; integrating the resulting intervals gives
\eqref{eq:J-no-root}. In the remaining case their first endpoint is
$z_1$ as above. Splitting there and at its reflected endpoint gives
\eqref{eq:J-root}. Thus the formulas retain the complete intervals,
including the shorter same-sign portion. The finite limits in these
displayed expressions give their claimed continuity.

Differentiating the complete expressions gives \eqref{eq:J-R-bounds};
the endpoint contributions cancel because the two channels agree at a
moving equality root. At $e=\pi/2$, $R=0$ for $s>v$, whereas its limit
at $e=0$ is $2(c\log c+v\log v-b\log b)$. These two endpoint formulas,
the sign of $R_e$, and the entropy formula for $s\leq v$ give the
displayed bound for $R$. The same differentiation gives the logarithmic
bound for $R_s$, whose singularity at $s=v$ is locally integrable.

For $s>v$, subtract $2s\log\sin e$ and integrate the resulting
separation derivative from the no-root threshold. It is
$-2\sqrt{(v^2-s^2\sin^2e)_+}/\sin e$, so the result is exactly $Z$.
Continuity includes $s=v$. If $s<v$, \eqref{eq:J-same-sign} gives
$J\leq0$, while
$\sqrt{v^2-s^2\sin^2t}\geq s\cos t$. Hence
$Z\geq-2s\log\sin e\geq E$. For $v=0$,
$J(0,s,e)=2s\log\sin e$ and $E=Z=0$ directly.
\end{proof}

\begin{lemma}[The unsplit root integral]\label{lem:root-majorant}
Let $a>0$, $0\leq v\leq a$, and
$I_4=\int_0^1\sqrt{1-t^4}\,dt$. Then
\begin{equation}\label{eq:F-root-bound}
 F(v,a):=\int_0^{\pi/2}Z(v,a\tan(e/2),e)\,de
       \leq\frac{4\sqrt2 I_4v^{3/2}}{\sqrt a},
 \qquad I_4<\frac{319}{360}.
\end{equation}
\end{lemma}
\begin{proof}
The case $v=0$ is immediate. Put $\tau=v/a\in(0,1]$. Tonelli's theorem
applies to the nonnegative integrand in \eqref{eq:root-Z} and retains
the triangular constraint $0<e<t<\pi/2$. At fixed $t$, substitute
$z=\tan(e/2)\sin t/\tau$. Since
$\tan(t/2)\sin t=1-\cos t$, this gives exactly
\[
 F=4v\tau\int_0^{\pi/2}
    \int_0^{\min(1,(1-\cos t)/\tau)}
       \frac{\sqrt{1-z^2}}{\sin^2t+\tau^2z^2}\,dz\,dt
   =4v\tau\int_0^1\sqrt{1-z^2}\,K(\tau z)\,dz,
\]
where, for $0<k\leq1$,
\[
 K(k)=\int_{\arccos(1-k)}^{\pi/2}\frac{dt}{\sin^2t+k^2}
 =\int_0^{(1-k)/\sqrt{k(2-k)}}\frac{dw}{1+k^2+k^2w^2}
 \leq\frac1{\sqrt{2k}}.
\]
The substitution in the second equality is $w=\cot t$. The last
inequality follows from $1+k^2\geq1$ and
$2(1-k)^2\leq2-k$. It also includes the collapsed interval at $k=1$.
The endpoint $k=0$ remains improper. Since
$\int_0^1\sqrt{1-z^2}\,z^{-1/2}dz=2I_4$, substitution proves the first
bound and simultaneously proves finiteness. Finally, for $0\leq u\leq1$,
$P(u)=1-u/2-u^2/8>0$ and
$P(u)^2-(1-u)=u^3/8+u^4/64\geq0$. Therefore
$I_4<1-1/10-1/72=319/360$.
\end{proof}

\begin{corollary}\label{cor:root-rational-constants}
For the actual phases $x,y\geq0$, $\rho=x+y\leq1$, let
$a=(\cos x+\cos y)/2$, $p=\sin x$, and $\kappa=4I_4/\pi^2$.
Then
\[
 \frac{\kappa p^{3/2}}{\sqrt a}\leq
 \begin{cases}(3/8)x^{3/2},&\rho\leq1/2,\\
               (5/12)x^{3/2},&\rho\leq1.
 \end{cases}
\]
\end{corollary}
\begin{proof}
On $\rho\leq1/2$, $a\geq1-\rho^2/4\geq15/16>(29/30)^2$.
With $\pi^2>197/20$ and Lemma~\ref{lem:root-majorant},
\[
 \frac\kappa{\sqrt a}<\frac{319}{87\pi^2}
               <\frac{6380}{17139}<\frac38,
\]
the final comparison being $51040<51417$. On $\rho\leq1$,
$c_e,c_d\geq7/8$, so $a\geq49/64$ and
\[
 \frac\kappa{\sqrt a}<\frac{1276}{315\pi^2}
               <\frac{25520}{62055}<\frac5{12},
\]
where $306240<310275$. Finally $0\leq p\leq x$.
\end{proof}

\begin{lemma}[The closed half-pi central decomposition]\label{lem:central-model}
Let $x,y\geq0$, $\rho=x+y\leq\pi/2$, and put
\begin{gather*}
 c_e=\cos(\rho/2),\quad c_d=\cos((x-y)/2),\quad
 a=c_ec_d,\\ p=\sin x,\quad q=\sin y,\quad
 v=p/2,\quad \delta=q/2.
\end{gather*}
For the endpoint finite part $E_\infty$ and scalar finite part $H$ of
the completed model in \eqref{scalar:completed-model}, its constant term
has the exact decomposition
\begin{equation}\label{eq:exact-central-decomposition}
 C_\infty=E_\infty+c_eH(x-y)+c_dH(x+y)-c\,a+E_c,
 \qquad c=\frac13+\frac4{\pi^2}(\log^2\pi-\log^22),
\end{equation}
where
\begin{equation}\label{eq:Ec-definition}
 E_c=\frac1{\pi^2}\int_0^{\pi/2}
 \left\{J(v,a\tan(t/2)+\delta,t)
       +J(v,|a\tan(t/2)-\delta|,t)
       -4a\tan(t/2)\log(\sin t)\right\}\,dt.
\end{equation}
The integral is continuous on the whole closed phase set. On that set
$a\geq1/2$, $v/a\leq1$, and
$b=2\arctan(q/(2a))\leq\rho\leq\pi/2$.
\end{lemma}
\begin{proof}
The phase identities give
\[
 2a-(p+q)=2c_d\{\cos(\rho/2)-\sin(\rho/2)\}\geq0,
 \qquad c_e,c_d\geq2^{-1/2}.
\]
The first inequality is strict in the interior $\rho<\pi/2$. In the
endpoint half of the defining completed finite-part integral, put
$T=2a\cot(t/2)$, $0<t\leq\pi/2$. Then $T>p+q$ and
\[
 (T-q)\sin t=2a(1+\cos t)-q\sin t\geq2a-q>p.
\]
Both generic channels are consequently in the no-root regime of
\eqref{eq:J-no-root}. Their sum is exactly
$4a\cot(t/2)\log\sin t$. In the central half, reflection
$t=\pi-z$ gives the two terms retained in \eqref{eq:Ec-definition}.
Subtracting this endpoint baseline before integration, the scalar pole
and finite-part terms are respectively
$c_eH(x-y)+c_dH(x+y)$ and $-ca$; the endpoint contribution is
$E_\infty$. This proves \eqref{eq:exact-central-decomposition} in the
strict interior, with the original finite-part normalization.

To justify the boundary, $0\leq v\leq1/2$,
$0\leq a\tan(t/2)\leq1$, and both channel parameters are at most
$3/2$. Lemma~\ref{lem:generic-cancellation} bounds the complete numerator
in \eqref{eq:Ec-definition} by
$10|\log\sin t|+4\log2$. This is integrable on $(0,\pi/2)$.
Dominated convergence proves continuity through every moving root,
zero coefficient, axis and boundary phase. The complete defining
finite-part integral for $C_\infty$ has the same integrable subtracted
endpoint envelopes; $E_\infty$ and $H$ are continuous for the same
reason. Scaling any boundary phase toward the origin and passing to
the limit proves the exact equality on the closed boundary. At its
possible degeneracy $T=p+q$ one has $t=\pi/2$, where $\log\sin t=0$;
the generic formulas have their common finite value. Finally,
$(\cos x+\cos y)\tan(\rho/2)=p+q\geq q$ proves $b\leq\rho$.
\end{proof}
